\documentclass[../../../include/open-logic-section]{subfiles}

\begin{document}

\olfileid[pt]{sth}{ord-arithmetic}{expo}
\olsection{Exponenciação Ordinal}

Passamos agora à exponenciação ordinal. Infelizmente, não há uma definição
\emph{simpática} sintética para a exponenciação ordinal.

Certo que \emph{existem} definições sintéticas explícitas. Eis uma.
Seja $\text{finfun}(\alpha,\beta)$ o conjunto de todas as funções $f
\colon \alpha \to \beta$ tais que $\Setabs{\gamma \in
\alpha}{f(\gamma) \neq 0}$ é equipotente a algum número natural.
Defina-se um bom ordenamento em $\text{finfun}(\alpha,\beta)$ por $f
\sqsubset g$ sse $f \neq g$ e $f(\gamma_0) < g(\gamma_0)$, onde
$\gamma_0 = \text{max}\Setabs{\gamma \in \alpha}{f(\gamma) \neq
g(\gamma)}$. Então pode definir-se $\ordexpo{\alpha}{\beta}$ como
$\ordtype{\text{finfun}(\alpha, \beta), \sqsubset}$. Potter usa
esta definição explícita e explica de seguida:
\begin{quote}
	A escolha desta ordenação determina-se unicamente pelo desejo de
	obter uma definição da exponenciação ordinal que obedeça à
	condição recursiva adequada\ldots, e é muito mais difícil de
	visualizar do que a soma ou o produto ordenados.
	\citep[p.~199]{Potter2004}
\end{quote}
Exatamente. Explicamos a adição como ``uma cópia de $\alpha$ seguida de uma
cópia de $\beta$'', e a multiplicação como ``uma $\beta$-sucessão de
cópias de $\alpha$''. Mas não temos nada de sucinto a dizer sobre
$\text{finfun}(\alpha, \gamma)$. Por isso, apresentaremos a
definição da exponenciação ordinal apenas \emph{por} recursão transfinita, O.s.:

\begin{defn}\ollabel{ordexporecursion}
\begin{align*}
	\ordexpo{\alpha}{0} &= 1\\
	\ordexpo{\alpha}{\beta\ordplus 1} &=\ordexpo{\alpha}{\beta} \ordtimes \alpha\\
	\ordexpo{\alpha}{\beta} &= \bigcup_{\delta < \beta}\ordexpo{\alpha}{\delta}& & \text{quando $\beta$ é um ordinal limite}
\end{align*}
\end{defn}

Se trabalhássemos \emph{como} teoristas de conjuntos, poderíamos querer explorar
algumas propriedades da exponenciação ordinal. Mas pouco mais temos a acrescentar,
salvo notar o fato pouco surpreendente de que a exponenciação ordinal não comuta.
Assim $\ordexpo{2}{\omega} =
\bigcup_{\delta < \omega}\ordexpo{2}{\delta} = \omega$, enquanto
$\ordexpo{\omega}{2} = \omega \ordtimes \omega$. Mas então não se deve
\emph{esperar} que comute, pois já não comuta nos naturais: $\ordexpo{2}{3} = 8 < 9 = \ordexpo{3}{2}$.

\begin{prob}
Usando indução transfinita, prove que, se se define
$\ordexpo{\alpha}{\beta} = \ordtype{\text{finfun}(\alpha, \beta),
\sqsubset}$, obtêm-se as equações recursivas de
\olref[sth][ord-arithmetic][expo]{ordexporecursion}.
\end{prob}

\end{document}
